%!TEX encoding = UTF-8 Unicode
%---------------------------------------------------------
\cvsection{科研与工作经历}

\begin{cventries}
  \cventry
    {博士后}
    {香港中文大学}
    {中国香港}
    {2025年6月 -- 至今}
    {
      \begin{cvitems}
        \item {开发 RNA 多组学分析方法与流程，研究疾病机制，探索潜在治疗靶点与生物标志物。}
      \end{cvitems}
    }

  \cventry
    {生物信息总监}
    {深圳安吉康尔科技有限公司}
    {中国广东深圳}
    {2018年5月 -- 2019年7月}
    {
      \begin{cvitems}
        \item {牵头开发自动化遗传病检测平台，负责系统架构设计、算法选型与分析流程实现。}
        \item {统筹一体化变异检测设备研发，推进软硬件集成与产品落地。}
      \end{cvitems}
    }

  \cventry
    {生物信息工程师}
    {欧蒙医学诊断（中国）有限公司}
    {中国北京}
    {2018年2月 -- 2018年4月}
    {
      \begin{cvitems}
        \item {参与耳聋检测的 NGS 分析模型开发与优化。}
      \end{cvitems}
    }

  \cventry
    {生物信息组长（血液肿瘤项目）}
    {北京橡鑫生物科技有限公司}
    {中国北京}
    {2017年3月 -- 2018年1月}
    {
      \begin{cvitems}
        \item {带领团队开发用于血液肿瘤的 NGS FLT3-ITD 检测方法，并与毛细管电泳结果进行定量比较。}
        \item {负责血液肿瘤体外诊断（IVD）产品的生物信息研发，建立分析流程，并为产品注册提供技术支持。}
      \end{cvitems}
    }

  \cventry
    {生物信息工程师（血液肿瘤项目）}
    {安诺优达基因科技（北京）有限公司}
    {中国北京}
    {2015年7月 -- 2017年2月}
    {
      \begin{cvitems}
        \item {负责血液肿瘤临床检测的生物信息分析，涵盖 NGS 数据处理与结果报告。}
        \item {独立开发分析流程自动运行与监控平台，支持 NGS 数据分析的无人值守运行。}
      \end{cvitems}
    }
\end{cventries}
